  \subsubsection{床関数で添字を二つに分割する}\label{節：ガウス型の添字の分割}
    ここからは,\bookterm{floor-function}{床関数}が\bookterm{sequence}{数列}の値ではなく,\textbf{どの項を参照するか}を指定する使い方を考える.
    正整数$n$に対して,
    \[
      \left\lfloor\frac{n+1}{2}\right\rfloor
      +\left\lfloor\frac n2\right\rfloor=n
    \]
    であり,二つの添字は$n$をできるだけ等しく分けた大きさになる.
    偶数では両方が同じ添字になり,奇数では隣り合う添字になることが手掛かりである.

    この形の問題では,まず偶奇で分けて\bookterm{floor-function}{床関数}を外し,それから\bookterm{difference}{階差}を取る方法がある.
    半分の添字との関係が得られれば,\,$1,\ 2\le n<4,\ 4\le n<8,\ldots$のように,
    2の累乗ごとの群で整理できる.
    定義の確認も必要で,例えば$n=1$で元と同じ添字を参照する式は,
    そのまま初項を決める式として使えないことがある.

    \noindent \bookblocklink{example-gauss-role-7}{problem}{answer}{【例題】}\par\nobreak
    \begin{enumerate}[label=(\arabic*),parsep=3pt,beginpenalty=10000,format=\bookitemlink{example-gauss-role-7}{problem}{answer}]
      \item \bookterm{sequence}{数列}を次のように定める.\bookterm{general}{一般項}を求めよ.
\BookMathLayout{            \[
              a_1=2,\qquad
              a_n=a_{\left\lfloor\frac{n+1}{2}\right\rfloor}
                  +a_{\left\lfloor\frac n2\right\rfloor}+n-1
              \quad(n\ge2).
            \]}{            \[
              \begin{aligned}
              a_1&=2,\\
              a_n&=a_{\left\lfloor\frac{n+1}{2}\right\rfloor}
                  +a_{\left\lfloor\frac n2\right\rfloor}+n-1
              \\              &\quad(n\ge2).
              \end{aligned}
            \]}
            必要ならば,\,$2^m\le n<2^{m+1}$を満たす整数$m\ge0$を用いてよい.
    \end{enumerate}

    \noindent \bookblocklink{example-gauss-role-7}{hint}{problem}{【方針】}\par\nobreak
    \begin{enumerate}[label=(\arabic*),parsep=3pt,beginpenalty=10000,format=\bookitemlink{example-gauss-role-7}{hint}{problem}]
      \item 偶奇で式を分け,\,$b_n=a_{n+1}-a_n$を考えます.
            \,$b_{2k}$と$b_{2k+1}$が$b_k+1$になることを示し,群ごとに足しましょう.
    \end{enumerate}

    \noindent \bookblocklink{example-gauss-role-7}{answer}{problem}{【解答】}\par\nobreak
    \begin{enumerate}[label=(\arabic*),parsep=3pt,beginpenalty=10000,format=\bookitemlink{example-gauss-role-7}{answer}{problem}]
      \item \,$2^m\le n<2^{m+1}$を満たす整数$m\ge0$に対して,
            \[
              a_n=(m+3)n-2^{m+1}+1.
            \]
    \end{enumerate}

    \noindent \bookblocklink{example-gauss-role-7}{solution}{problem}{【解法】}\par\nobreak
    \begin{enumerate}[label=(\arabic*),parsep=3pt,beginpenalty=10000,format=\bookitemlink{example-gauss-role-7}{solution}{problem}]
      \item \,$n\ge2$では二つの添字は1以上$n-1$以下なので,\bookterm{sequence}{数列}は順に定義できる.
            偶奇で分けると,\,$k\ge1$について,
            \[
              \begin{aligned}
                a_{2k}&=2a_k+2k-1,\\
                a_{2k+1}&=a_{k+1}+a_k+2k.
              \end{aligned}
            \]
            \,$b_n=a_{n+1}-a_n$とおくと,両式の差から$b_{2k}=b_k+1$.
            また,\,$a_{2k+2}=2a_{k+1}+2k+1$を使うと$b_{2k+1}=b_k+1$である.
            \,$a_2=5$より$b_1=3$なので,
            \[
              b_{2k}=b_{2k+1}=b_k+1,\qquad b_1=3.
            \]

            \,$2^m\le n<2^{m+1}$の範囲では$b_n=m+3$となる.
            \,$m=0$では$b_1=3$と一致する.
            次の群の添字を半分にして床を取ると前の群に入り,値が1増えるから,
            群番号$m$についての帰納法でこの式が示される.

            従って,完全な群と最後の群を分けて足すと,
            \begin{align*}
              a_n={}&2+\sum_{j=0}^{m-1}2^j(j+3)\\
                   &+(n-2^m)(m+3).
            \end{align*}
            ここで,
            \[
              \sum_{j=0}^{m-1}2^j=2^m-1,
            \]
            \[
              \sum_{j=0}^{m-1}j2^j=(m-2)2^m+2
            \]
            を使う.後者は,和を$T$とおき$2T-T$を計算すれば得られる.
            \,$m=0$でも空の和として両式は成り立つ.
            これを代入すると,
            \[
              a_n=(m+3)n-2^{m+1}+1
            \]
            を得る.添字を分割する問題から,\bookterm{difference}{階差}を通じて群\bookterm{sequence}{数列}へ戻ることができた.
    \end{enumerate}
